\chapter{Introduzione}
Negli ultimi anni, la Realtà Aumentata (AR) ha completato la sua transizione da tecnologia sperimentale a strumento pervasivo, capace di ridefinire l'interazione uomo-macchina in settori che spaziano dall'industria 4.0 alla medicina, fino all'istruzione e all'intrattenimento. A differenza della Realtà Virtuale, che isola l'utente in un ambiente digitale, l'AR arricchisce la percezione del mondo fisico sovrapponendovi elementi virtuali in tempo reale e in modo geometricamente coerente.Nonostante il potenziale rivoluzionario, la diffusione di massa delle applicazioni AR ha dovuto affrontare una barriera significativa: la frammentazione hardware e software. Per lungo tempo, lo sviluppo di soluzioni AR ha richiesto l'utilizzo di SDK(Software Development Kit) nativi differenti, come Apple ARKit per i dispositivi iOS e Google ARCore per l'ecosistema Android. Questa dicotomia ha storicamente comportato un aumento dei costi di produzione, problemi di portabilità del codice e una complessa gestione degli aggiornamenti.In questo panorama si colloca AR Foundation, il framework cross-platform sviluppato da Unity. AR Foundation si propone come uno strato di astrazione software che unifica le API dei principali sistemi operativi, permettendo agli sviluppatori di progettare un'architettura logica universale (“write once, run anywhere”). Attraverso l'uso di componenti modulari (come i Manager per il tracciamento delle superfici o dei volti), il framework non solo ottimizza il flusso di lavoro, ma garantisce l'accesso alle funzionalità hardware più avanzate dei dispositivi moderni senza perdere in prestazioni.
In questo lavoro, risultato dell'attivitá progettuale svolta presso il dipartimento di Ingegneria "Enzo Ferrari" dell'Universitá di Modena e Reggio Emilia, l'obiettivo è analizzare l'architettura e le potenzialità di AR Foundation come soluzione al problema della frammentazione nello sviluppo di applicazioni in Realtà Aumentata. Nella prima parte della tesi, a carattere compilativo, viene fornito un quadro comparativo dettagliato tra AR Foundation e gli SDK nativi sottostanti ARCore e ARKit, valutando l'impatto dell'astrazione in termini di funzionalità supportate, flessibilità architetturale e prestazioni. Successivamente, la trattazione si focalizzerà su una fase sperimentale, incentrata sulla progettazione e sull'implementazione di un'applicazione dimostrativa, al fine di quantificare sul campo l'efficienza del framework nei principali casi d'uso industriali o applicativi.